\section{The Denavit--Hartenberg frame convention}
\label{sec:l10-dh}
The product in Equation~\eqref{eq:l10-product} reduces the direct-kinematics
problem to finding the pose of each link relative to its predecessor.
Arbitrary link frames are valid, but they can make those local matrices
unnecessarily complicated. The \emph{Denavit--Hartenberg} (DH) convention
chooses frames so that each local transformation has just four parameters:
two translations and two rotations. For a one-degree-of-freedom joint,
only one of these parameters varies with joint motion.

We use the \emph{standard DH convention}. In particular, joint $i$ acts
about or along $z_{i-1}$, and the local matrix maps coordinates from frame
$i$ to frame $i-1$. Other conventions, including modified DH, use different
frame assignments and matrix formulas. A parameter table must always be
used with the convention that defines it.

\subsection{Joint axes and common normals}
Let $J_i$ denote the axis of joint $i$: its rotation axis for a revolute
joint or its translation direction, represented by a chosen axis line,
for a prismatic joint. For the generic construction, consider neighboring
axes $J_i$ and $J_{i+1}$ that are \emph{skew}: they neither intersect nor
are parallel.

The shortest segment joining two skew lines is perpendicular to both.
The line containing that segment is their \emph{common normal}, denoted
$\mathrm{CN}_i$. It intersects the axes at their closest points. Extending
either joint axis farther does not change these points or the separation.
The common normal is a geometric construction; it need not follow the
shape of the physical link joining the joints.

\subsection{Assigning an interior link frame}
For $i=1,\ldots,n-1$, place frame $i$ on link $i$ using the following rules:
\begin{enumerate}
\item Put $O_i$ at the intersection of $J_{i+1}$ and $\mathrm{CN}_i$.
\item Align $z_i$ with $J_{i+1}$ and choose its positive direction.
\item Direct $x_i$ along $\mathrm{CN}_i$, from $J_i$ toward $J_{i+1}$.
      The axis starts at $O_i$ and points in that same direction.
\item Complete a right-handed frame with
      $\hat y_i=\hat z_i\times\hat x_i$.
\end{enumerate}
Here hats denote unit direction vectors. The common-normal construction
makes $x_i$ perpendicular to $z_i$. Consequently, their directions determine
$y_i$ uniquely; a drawing can omit $y_i$ without leaving the frame undefined.

The index shift is essential: $z_i$ follows joint $i+1$, whereas joint $i$
connects links $i-1$ and $i$ and acts about or along $z_{i-1}$. There are
$n$ joints, $n$ moving links plus the base, and $n+1$ link frames numbered
$0,\ldots,n$. An interior link $i$ lies between joints $i$ and $i+1$;
the terminal link has no physical joint $n+1$.

Either orientation of a joint axis can be chosen initially. Reversing it
changes the associated signs and frame assignments, so all parameters and
joint-coordinate definitions must be made consistent with that choice.
Once chosen, the frames remain attached to their links as the arm moves.

\subsection{End frames and nonunique constructions}
The interior rule cannot be applied literally at both ends of the chain.
For frame 0, choose $O_0$ on $J_1$, align $z_0$ with $J_1$, and choose a
perpendicular $x_0$. This base frame need not coincide with the world frame.
At the other end, choose a terminal frame compatible with the DH
construction; its $z_n$ is an auxiliary axis, not an additional joint.
A desired tool frame can be connected to it by a fixed transform as in
Equation~\eqref{eq:l10-offsets}.

The skew-axis rule also needs choices in special geometries. Distinct
parallel axes have infinitely many common normals. Intersecting,
nonparallel axes have zero separation, with a normal direction
perpendicular to both. Coincident axes leave still more freedom.
These situations are common in robot designs. The transformation derived
below still applies once a consistent DH frame assignment has been made;
the frame assignment need not be unique.

\section{Four parameters for one adjacent-frame transformation}
\label{sec:l10-dh-parameters}
To relate frames $i-1$ and $i$, introduce a temporary frame $i'$.
Its origin is where $J_i$ meets $\mathrm{CN}_i$. Choose
$\hat z_{i'}=\hat z_{i-1}$ and $\hat x_{i'}=\hat x_i$, with $y_{i'}$
completing the right-handed frame. This frame combines the incoming joint
axis with the outgoing common-normal direction.

\begin{figure}[htbp]
\centering\begin{tikzpicture}[x=1cm,y=1cm]
\definecolor{dhgreen}{RGB}{25,145,65}
\definecolor{dhorange}{RGB}{235,125,20}
% Layout follows the board: two shaped links above their axis constructions.
\coordinate (L) at (0.8,3.7);
\coordinate (M) at (6,4.4);
\coordinate (N) at (8.8,6.5);
\coordinate (A) at (5.5,.5);
\coordinate (P) at (5.8,2.84);
\coordinate (B) at (6.9,2.02);
\coordinate (C) at (1.25,1.9);
% Axis lines behind the solid links.
\draw[red!75!black,thick] (.42,5.15)--(1.65,.38);
\draw[red!75!black,thick] (5.42,-.12)--(6.22,6.12);
\draw[red!75!black,thick] (6.62,1.35)--(9.16,7.34);
\node[red!75!black,above] at (.42,5.15) {$J_{i-1}$};
\node[red!75!black,above] at (6.22,6.12) {$J_i$};
\node[red!75!black,above] at (9.16,7.34) {$J_{i+1}$};
% Far link and its small visible lower edge.
\path[draw,thick,fill=gray!18] (5.56,4.26)
 .. controls (5.9,4.02) and (6.28,4.15) .. (6.68,4.75)
 .. controls (7.65,6.0) and (8.22,5.83) .. (8.96,6.21)
 .. controls (9.47,6.5) and (9.24,7.0) .. (8.8,6.98)
 .. controls (8.15,6.97) and (7.55,6.26) .. (6.96,5.61)
 .. controls (6.35,4.94) and (6.13,4.83) .. (5.7,4.74) -- cycle;
\path[draw,thick,fill=white] (5.56,4.4)
 .. controls (5.9,4.16) and (6.28,4.29) .. (6.68,4.89)
 .. controls (7.65,6.14) and (8.22,5.97) .. (8.96,6.35)
 .. controls (9.47,6.64) and (9.24,7.14) .. (8.8,7.12)
 .. controls (8.15,7.11) and (7.55,6.4) .. (6.96,5.75)
 .. controls (6.35,5.08) and (6.13,4.97) .. (5.7,4.88) -- cycle;
% Near link, curved like the board drawing rather than a stick diagram.
\path[draw,thick,fill=gray!18] (.35,3.34)
 .. controls (.12,3.75) and (.62,4.14) .. (1.12,3.9)
 -- (1.5,3.65)
 .. controls (3.3,3.96) and (4.75,3.86) .. (5.55,4.48)
 .. controls (6.18,4.85) and (6.65,4.53) .. (6.52,4.08)
 .. controls (6.4,3.73) and (5.92,3.77) .. (5.37,3.8)
 .. controls (3.91,3.96) and (2.48,3.57) .. (1.35,3.28)
 .. controls (.87,3.16) and (.55,3.19) .. (.35,3.34);
\path[draw,thick,fill=white] (.35,3.48)
 .. controls (.12,3.89) and (.62,4.28) .. (1.12,4.04)
 -- (1.5,3.79)
 .. controls (3.3,4.1) and (4.75,4.0) .. (5.55,4.62)
 .. controls (6.18,4.99) and (6.65,4.67) .. (6.52,4.22)
 .. controls (6.4,3.87) and (5.92,3.91) .. (5.37,3.94)
 .. controls (3.91,4.1) and (2.48,3.71) .. (1.35,3.42)
 .. controls (.87,3.3) and (.55,3.33) .. (.35,3.48);
\node[rotate=6] at (3.2,4.25) {Link $i-1$};
\node[rotate=45] at (7.45,5.91) {Link $i$};
\foreach \p in {L,M,N} {\draw[fill=white,thick] (\p) ellipse (.09 and .055);}
% Common normals and the three frames below the links.
\draw[sensorframe,very thick] (C)--node[below,sloped] {$\mathrm{CN}_{i-1}$}(A);
\draw[sensorframe,very thick] (P)--node[below,sloped] {$\mathrm{CN}_i$}(B);
\draw[axis] (A)--++(.12,.94) node[right] {$z_{i-1}$};
\draw[axis] (A)--++(1.1,-.36) node[right] {$x_{i-1}$};
\draw[axis,dhorange] (P)--++(.10,.78) node[left] {$z_{i'}$};
\draw[axis,dhorange] (P)--++(.66,-.492);
\node[dhorange] at (6.23,2.69) {$x_{i'}$};
\draw[axis] (B)--++(.44,1.04) node[right] {$z_i$};
\draw[axis] (B)--++(.86,-.64) node[right] {$x_i$};
\draw[axis] (B)--++(.87,.36) node[right] {$y_i$};
\foreach \p in {A,P,B} {\fill (\p) circle (.035);}
\node[below left] at (A) {$O_{i-1}$};
\node[left,dhorange] at (P) {$O_{i'}$};
\node[below left] at (B) {$O_i$};
% Dimension offsets preserve the same directions as their measured segments.
\draw[projection,dhgreen] (A)--(4.65,.5);
\draw[projection,dhgreen] (P)--(4.95,2.84);
\draw[<->,dhgreen,thick] (4.65,.5)--node[left] {$d_i$}(4.95,2.84);
\draw[projection,dhgreen] (P)--(5.6,2.35);
\draw[projection,dhgreen] (B)--(6.7,1.53);
\draw[<->,dhgreen,thick] (5.6,2.35)--node[below,sloped] {$a_i$}(6.7,1.53);
% Angle guides compare directions translated to the same origin.
\draw[projection] (P)--++(1.12,-.369);
\draw[->,dhgreen,thick] (6.53,2.60) .. controls (6.57,2.49) and (6.52,2.38) .. (6.43,2.37);
\node[dhgreen] at (6.65,2.98) {$\theta_i$};
\draw[projection] (B)--++(.17,1.34);
\draw[->,dhgreen,thick] (7.05,3.19) .. controls (7.16,3.22) and (7.28,3.17) .. (7.36,3.10);
\node[dhgreen,above] at (7.23,3.32) {$\alpha_i$};
\end{tikzpicture}

\caption{DH construction for two neighboring links. Red lines extend the
three joint axes; blue lines show their common normals. The intermediate
frame $i'$ is orange, and the four parameters are marked in green.
Common normals are perpendicular to their joint axes in space, even when
the perspective projection does not show a right angle. Omitted $y$ axes
follow from the right-hand rule.}
\label{fig:l10-dh-geometry}
\end{figure}

The four DH parameters describe two successive changes of frame:
\begin{itemize}
\item $d_i$ is the signed displacement from $O_{i-1}$ to $O_{i'}$
      along $z_{i-1}$.
\item $\theta_i$ is the angle from $x_{i-1}$ to $x_{i'}=x_i$
      about $z_{i-1}$.
\item $a_i$ is the displacement from $O_{i'}$ to $O_i$ along $x_i$.
\item $\alpha_i$ is the angle from $z_{i'}=z_{i-1}$ to $z_i$
      about $x_i$.
\end{itemize}
Angles follow the right-hand rule about their named positive axes.
A positive angle appears counterclockwise when viewed from the positive
end of that axis toward the origin. This viewing direction matters in a
perspective drawing. The lengths $a_i,d_i$ have units of length; the
angles $\alpha_i,\theta_i$ are measured in radians unless stated otherwise.

With $x_i$ directed from one skew axis toward the next, $a_i$ is their
nonnegative separation. The offset $d_i$ may be positive, zero, or negative:
it is a directed displacement, not an unsigned distance.

\subsection{Reading the parameters without the temporary frame}
The temporary frame helps derive the matrix, but it need not be drawn for
every joint. Since the two displacement directions are perpendicular,
the vector between the link-frame origins satisfies
\begin{equation}
 \overrightarrow{O_{i-1}O_i}=d_i\hat z_{i-1}+a_i\hat x_i.
 \label{eq:l10-dh-displacement}
\end{equation}
All vectors in this equation must be expressed in the same reference
frame. Taking projections gives
\[
 d_i=\hat z_{i-1}\cdot\overrightarrow{O_{i-1}O_i},\qquad
 a_i=\hat x_i\cdot\overrightarrow{O_{i-1}O_i}.
\]
Thus the phrase ``distance between the origins along an axis'' means a
\emph{signed projection}, not the length of the diagonal joining them.
For example, if $a_i=0.4\,\mathrm m$ and $d_i=0.3\,\mathrm m$, the
origins are $0.5\,\mathrm m$ apart, but neither DH translation equals
$0.5\,\mathrm m$.

Likewise, $\theta_i$ can be read directly from $x_{i-1}$ to $x_i$ about
$z_{i-1}$, and $\alpha_i$ from $z_{i-1}$ to $z_i$ about $x_i$.
Directions can be translated to a common point to compare their angles.
The auxiliary frame adds no joint, degree of freedom, or extra row to the
DH table.

\section{Deriving the standard DH matrix}
\label{sec:l10-dh-matrix}
The temporary frame splits the desired transformation into
\begin{equation}
 H_i^{i-1}=H_{i'}^{i-1}H_i^{i'}.
 \label{eq:l10-dh-factor}
\end{equation}
From frame $i-1$ to frame $i'$, the origin moves by $d_i$ along
$z_{i-1}$ and the axes rotate by $\theta_i$ about that same axis:
\begin{equation}
 H_{i'}^{i-1}=\begin{bmatrix}
 \cos\theta_i&-\sin\theta_i&0&0\\
 \sin\theta_i& \cos\theta_i&0&0\\
 0&0&1&d_i\\0&0&0&1
 \end{bmatrix}.
 \label{eq:l10-dh-z}
\end{equation}
From $i'$ to $i$, the origin moves by $a_i$ along $x_{i'}=x_i$ and
the axes rotate by $\alpha_i$ about that direction:
\begin{equation}
 H_i^{i'}=\begin{bmatrix}
 1&0&0&a_i\\
 0&\cos\alpha_i&-\sin\alpha_i&0\\
 0&\sin\alpha_i& \cos\alpha_i&0\\
 0&0&0&1
 \end{bmatrix}.
 \label{eq:l10-dh-x}
\end{equation}
The translation in this second matrix is along $x$, even though the
twist changes the direction of $z$. Matrix multiplication gives
\begin{equation}
 \boxed{H_i^{i-1}=\begin{bmatrix}
 c_{\theta_i}&-s_{\theta_i}c_{\alpha_i}&s_{\theta_i}s_{\alpha_i}&a_ic_{\theta_i}\\
 s_{\theta_i}&c_{\theta_i}c_{\alpha_i}&-c_{\theta_i}s_{\alpha_i}&a_is_{\theta_i}\\
 0&s_{\alpha_i}&c_{\alpha_i}&d_i\\
 0&0&0&1
 \end{bmatrix},}
 \label{eq:l10-dh-full}
\end{equation}
where $c_\beta=\cos\beta$ and $s_\beta=\sin\beta$.
The translation column is expressed in frame $i-1$: the displacement
$a_i$ is resolved along its $x$ and $y$ directions by $\theta_i$.

Using $4\times4$ elementary rotation and translation matrices, the same
factorization is
\[
 H_i^{i-1}=\operatorname{Rot}_z(\theta_i)
 \operatorname{Trans}_z(d_i)
 \operatorname{Trans}_x(a_i)\operatorname{Rot}_x(\alpha_i).
\]
This describes successive local frame placements: after the $z$ rotation,
the $x$ operations use the rotated $x$ direction. When the matrix acts on
a coordinate column, its rightmost factor acts first, as usual.
A translation along an axis commutes with a rotation about that same
axis, but the $z$ and $x$ groups cannot generally be interchanged.

\subsection{Where the joint variable enters}
The parameters $a_i$ and $\alpha_i$ describe fixed link geometry.
For a revolute joint, $d_i$ is also constant and $\theta_i$ varies.
For a prismatic joint, $\theta_i$ is constant and $d_i$ varies:
\begin{center}
\begin{tabular}{lll}
\toprule
Joint type & Variable parameter & Constant parameters\\
\midrule
Revolute & $\theta_i=q_i+\theta_{i,0}$ & $a_i,\alpha_i,d_i$\\
Prismatic & $d_i=q_i+d_{i,0}$ & $a_i,\alpha_i,\theta_i$\\
\bottomrule
\end{tabular}
\end{center}
The constants $\theta_{i,0}$ and $d_{i,0}$ allow the joint coordinate's
zero to differ from the DH parameter's zero. If the zeros coincide,
these reduce to $\theta_i=q_i$ and $d_i=q_i$, respectively. These formulas
assume $q_i$ increases in the chosen positive DH direction; an oppositely
signed encoder coordinate must be converted accordingly.

This explains why four parameters still give a one-variable local map
$H_i^{i-1}(q_i)$. The constants can be obtained from dimensions, a CAD
model, or calibration. Moving another joint changes the world pose of
these frames but does not introduce another independent variable into
their relative transformation.

\subsection{Worked example: checking one DH row}
Consider a revolute joint with zero angular offset and
\[
 a_i=0.4\,\mathrm m,\qquad d_i=0.3\,\mathrm m,\qquad
 \alpha_i=\pi/2,\qquad q_i=\theta_i=\pi/2.
\]
Substitution into Equation~\eqref{eq:l10-dh-full} gives
\[
 H_i^{i-1}=\begin{bmatrix}
 0&0&1&0\\1&0&0&0.4\\0&1&0&0.3\\0&0&0&1
 \end{bmatrix}.
\]
The origin is at $(0,0.4,0.3)\,\mathrm m$ in frame $i-1$.
The first rotation column shows that $x_i$ points along $y_{i-1}$;
the third shows that $z_i$ points along $x_{i-1}$. The columns are
orthonormal and right-handed. These geometric checks are often easier
than finding a sign error by rereading the entire formula.

If the same constant values $a_i=0.4\,\mathrm m$, $\alpha_i=\pi/2$,
and $\theta_i=\pi/2$ instead describe a prismatic joint with $d_i=q_i$,
only the third translation entry changes as that joint moves. A DH row
must therefore specify which parameter is the joint variable, not merely
four numbers at one configuration.

\clearpage
\section{From a robot drawing to direct kinematics}
\label{sec:l10-dh-procedure}
The DH construction turns the geometric problem into a repeatable sequence:
\begin{enumerate}
\item Identify and number the $n$ physical joint axes $J_1,\ldots,J_n$.
      Record each joint's type and positive motion direction.
\item Construct the $n-1$ common normals between neighboring joint axes.
      Resolve nonunique choices consistently.
\item Assign frames $0,\ldots,n$, including the base and terminal choices.
      Drawing $x$ and $z$ is sufficient when right-handedness is understood.
\item Make a DH table with one row per joint, $i=1,\ldots,n$.
      Record $a_i,\alpha_i,d_i,\theta_i$ and identify the variable entry.
\item Substitute each row into Equation~\eqref{eq:l10-dh-full} to obtain
      $H_i^{i-1}(q_i)$.
\item Multiply $H_n^0=H_1^0H_2^1\cdots H_n^{n-1}$. Include the base
      and tool offsets when the requested frames differ from 0 and $n$.
\end{enumerate}
There is no row 0 in this indexing. There are $n$ matrices even though
only $n-1$ common normals join physical joint pairs: the terminal-frame
choice supplies the remaining geometry. The temporary frames $i'$ are
derivation aids and do not change these counts.

Before expanding a long product, check one simple configuration. Verify
that each origin lies where the drawing predicts, that the rotation
columns describe the intended axes, and that a positive change in $q_i$
moves the model in the intended direction. Different consistent frame
choices may produce different DH tables. They should give the same
physical tool pose when expressed in the same world and tool frames,
with the corresponding joint-coordinate definitions.

The planar RR matrices derived earlier are a useful consistency check.
Setting $a_1=\ell_1$, $a_2=\ell_2$, $\alpha_1=\alpha_2=0$,
$d_1=d_2=0$, and $\theta_i=q_i$ in the spatial DH formula reproduces
those planar transformations, with an additional unchanged $z$ coordinate.
